\section{TPRS: Candidate Generation with Learned Joint Selection}

The diagnosis in Section~2 establishes that paralog alignment fails for two independent reasons: (1) the human query protein is too diverged from the target species' paralog (cross-species identity ceiling, 0.22--0.30 F1 loss), and (2) tools cannot distinguish the correct paralog copy among similar candidates (locus confusion, 0.22--0.38 F1 loss). TPRS is an alignment method designed to address both failure modes. It operates in two stages (Fig.~\ref{fig:tprs_pipeline}): Stage~1 generates an expanded candidate pool along three dimensions where single-tool, single-query approaches are vulnerable; Stage~2 performs learned joint selection over all (locus, query-type) pairs.

Concretely, let $Q$ be a human query protein and $G$ the target genome. Stage~1 constructs a candidate set
\begin{equation}
\mathcal{C} = \{(L, q) \mid L \in \mathcal{L},\; q \in \{Q, P\}\},
\label{eq:candidate_set}
\end{equation}
where $\mathcal{L}$ is the set of candidate genomic loci and $P$ denotes the target species' own protein (retrieved from the canonical proteome via CDS annotation intersection). Stage~2 defines a learned scoring function $f_\theta: \mathbb{R}^{25} \to [0,1]$ and selects the optimal pair:
\begin{equation}
(L^*, q^*) = \arg\max_{(L, q) \in \mathcal{C}} f_\theta(\mathbf{x}_{L, q}),
\label{eq:selection}
\end{equation}
returning the corresponding miniprot spliced alignment as the final output.

\subsection{Design Logic: From Diagnosis to Algorithm}

Each root cause maps directly to a specific architectural decision:

\begin{itemize}
    \item \textbf{RC1 (identity ceiling) $\rightarrow$ query-identity dimension.} The forced-locus experiment shows that no C-stage algorithm can produce accurate spliced alignments from a diverged human query. The candidate pool must therefore include query proteins with high sequence identity to the target. The target species' own protein achieves $\sim$100\% amino acid identity, eliminating the identity barrier entirely. Rather than treating this as a separate fix, TPRS incorporates the species-matched protein as an additional \emph{query type} in the candidate pool---the router may select either the human query or the species protein depending on which produces the better alignment at a given locus.

    \item \textbf{RC2 (locus confusion) $\rightarrow$ tool-diversity dimension + learned selection.} The correct-locus rate collapse shows that no single tool reliably finds the correct paralog locus. The candidate pool must therefore include hits from diverse search strategies. Furthermore, no fixed heuristic can reliably rank candidates, because the discriminating signal involves non-linear interactions between alignment quality, genomic context, and query identity. A learned model is required.
\end{itemize}

This mapping yields a two-stage alignment method: \textbf{Candidate Generation} (Phases 1--3, expanding the candidate space) and \textbf{Learned Joint Selection} (Phase~4, scoring and selecting among candidates). Critically, TPRS does \emph{not} independently fix RC1 then RC2. It performs a single operation---joint selection over an expanded candidate pool---that implicitly addresses both root causes because the correct (locus, query-type) pair is made available and then identified through learned scoring.

\subsection{Stage 1: Candidate Generation}

Candidate generation expands the search space along three orthogonal dimensions to ensure the correct (locus, query-type) pair enters the candidate pool. This stage comprises Phases 1--3, corresponding to the three dimensions.

\subsubsection{Dimension 1: Tool Diversity (Phase~1)}

The first dimension runs four complementary A-stage (search/alignment) tools with relaxed parameter settings to maximize recall, using the human query protein against the target genome: \texttt{diamond}~\cite{buchfink2021diamond} (k-mer based translated search), \texttt{last}~\cite{kielbasa2011last} (adaptive seed scoring), \texttt{miniprot}~\cite{li2023miniprot} in ultra-sensitive mode (minimizer-based spliced alignment), and \texttt{mmseqs2}~\cite{steinegger2017mmseqs2} (k-mer sensitivity search). Tool selection is motivated by Exp~3: standard (default-parameter) versions of these four tools exhibit distinct species-level recall patterns (Table~\ref{tab:s_exp3_recall}), their 4-tool union recall@50 (0.53--0.81 across species) substantially exceeds any single tool (best: 0.20--0.53), and best 2-tool oracle combinations vary by species---no single pair dominates. Ablation analysis (Section~4.2, E6) confirms their complementarity: removing any single tool reduces best-overlap F1 by $\Delta = -0.31$ to $-0.33$, and no tool dominates---each removal causes a similar-magnitude drop.

Raw hits from the four tools are merged using a chromosome- and strand-aware clustering procedure: within each (chromosome, strand) group, hits are sorted by start coordinate and merged if the distance between adjacent groups is $\leq$50~kb. Each merged cluster is then expanded by $\pm$50~kb to capture flanking exons and regulatory splice signals:
\begin{equation}
\mathcal{L}_{1} = \text{Cluster}\!\left(\bigcup_{t \in \mathcal{T}} t(Q, G),\; \delta\right), \quad \delta = 50\,\text{kb},
\label{eq:clustering}
\end{equation}
where the 50~kb padding exceeds the upper quartile of intron lengths in vertebrate genomes ($\sim$40~kb), ensuring that candidate regions fully contain gene bodies. The merged and padded regions form the candidate locus set for each query protein.

\subsubsection{Dimension 2: Query Identity (Phase~3)}

The second dimension addresses RC1 directly. For each candidate locus from Phase~1, we identify the best-matching gene annotation by intersecting the locus with the target species' CDS annotations, and retrieve the corresponding protein sequence from the Ensembl canonical proteome. This use of CDS annotations is not a shortcut but a direct consequence of the causal diagnosis: the forced-locus experiment (Section~2.3) establishes that no C-stage algorithm can exceed $\sim$0.71 exon F1 on paralogs with a cross-species query---the human query protein simply lacks the necessary sequence motifs, and the bottleneck is the query itself, not the alignment engine. The species-matched protein, sharing $\sim$100\% amino acid identity with the target, is the only mechanism that breaks this ceiling, and its retrieval requires precisely the annotation resources that model-organism genomes already provide.

Miniprot is then run \emph{twice} per candidate locus---once with the human query protein and once with the species-matched protein:
\begin{equation}
A_{L, q} = \text{miniprot}(q, G[L]), \quad q \in \{Q, P\},
\label{eq:dual_query}
\end{equation}
using \texttt{-N~3} (limiting secondary alignments), single-threaded for reproducibility. When default parameters produce no spliced alignment at a locus, we retry with zero gap-open penalty (\texttt{-p~0.0}) as a fallback; fewer than 2\% of candidate loci require this fallback. Because the species protein shares $\sim$100\% amino acid identity with the target gene product, the cross-species identity ceiling is eliminated for this query type. The forced-locus experiment (Table~\ref{tab:forced_locus}) confirms this mechanism: with the human query at the correct paralog locus, miniprot F1 caps at 0.64--0.71, but with the species protein, the identity barrier is removed and alignment quality recovers to near-1:1 levels (0.92--0.98). The only change between the human-query and species-protein runs is the query protein identity.

From each miniprot spliced alignment output, we extract eight C-stage features:
\begin{itemize}
    \item \texttt{c\_identity}: nucleotide-level alignment identity (matches / aligned length). For species-matched protein queries, this typically exceeds 0.95.
    \item \texttt{c\_coverage}: fraction of the query protein length covered by the predicted CDS ($\text{CDS\_length} / (3 \times \text{query\_length})$).
    \item \texttt{c\_n\_exons}: number of predicted exons in the spliced alignment.
    \item \texttt{c\_cds\_length}: total predicted CDS length in base pairs.
    \item \texttt{c\_composite\_score}: a combined quality score:
\begin{equation}
c_{\mathrm{comp}} = c_{\mathrm{id}} \cdot \min(1,\,c_{\mathrm{cov}}) \cdot \min\!\left(1,\,\frac{c_{\mathrm{exons}}}{3}\right),
\label{eq:composite}
\end{equation}
rewarding high identity, full coverage, and multi-exon gene structure simultaneously.
    \item \texttt{c\_has\_hit}: binary indicator of whether miniprot produced any spliced alignment at this locus.
    \item \texttt{c\_hit\_score}: raw miniprot alignment score of the top hit.
    \item \texttt{c\_rank}: rank of this alignment among all miniprot hits at this locus.
\end{itemize}

\subsubsection{Dimension 3: Prior Knowledge (Phase~2)}

Proteins whose orthologous group members are not found by any A-stage tool at the primary locus require an additional recovery mechanism. We identify blind-miss proteins (no tool hit at the primary locus within K=200 candidates) and recover them using Ensembl Compara orthology annotations~\cite{herrero2016ensembl} combined with protein-level search. For each blind-miss protein, we first identify paralog loci where A-stage tools produced hits; for these loci, we retrieve the corresponding target-species protein from the canonical proteome (using Ensembl gene IDs from CDS annotations), then run MMseqs2 easy-search~\cite{steinegger2017mmseqs2} with this species protein against the target genome. If MMseqs2 hits overlap the primary ortholog locus, the protein is considered recovered, and the genomic coordinates of the co-ortholog gene are added as a candidate region.

Ensembl Compara is a curated, experimentally validated orthology database with no learned parameters, so this step does not introduce overfitting risk. The recovery is applied to the full query set before the train/test split, ensuring that every protein in the test set has the same opportunity for locus recovery as those in the training set.

The final candidate locus set is the union of Phase~1 tool-discovered and Phase~2 recovered loci:
\begin{equation}
\mathcal{L} = \mathcal{L}_{1} \cup \mathcal{L}_{2}.
\label{eq:locus_union}
\end{equation}
This phase recovers 733 of 735 proteins (99.7\%) that would otherwise have no candidate at the primary locus. The two unrecovered proteins correspond to cases where Ensembl Compara annotates no co-ortholog for the human query in that species. Phase~2 is essential for the affected subset ($\sim$27\% of queries): when co-ortholog recovery is removed, the affected proteins drop to a mean best-overlap F1 of 0.21 (vs.\ 0.64 with recovery, Section~4.2, E4).

\subsubsection{Candidate Pool Output}

The output of Stage~1 is a set of (locus, query-type) pairs per protein. With 4 tools generating hits, dual query types per locus, and co-ortholog recovery for missed proteins, the candidate pool averages 9.2 (mouse) to 16.8 (zebrafish) pairs per protein. Each pair has an associated miniprot spliced alignment and extracted feature vector.

\subsection{Stage 2: Learned Joint Selection}

The second stage addresses RC2 by training a learned router to select the best (locus, query-type) pair for each protein.

\subsubsection{Problem Formulation}

For each query protein $p$, let $\mathcal{C}_p = \{(L_1, t_1), \dots, (L_k, t_k)\}$ be the set of candidate pairs after Stage~1. The routing task is: given feature vectors $\mathbf{x}_i$ for each candidate pair, select the pair whose miniprot spliced alignment maximizes exon F1.

We formulate this as a regression problem rather than a classification problem for two reasons. First, the number of candidates per protein varies, making a fixed-size classifier output difficult. Second, regressing exon F1 directly provides a continuous quality score that naturally ranks candidates within each protein. At inference time, the predicted F1 scores for all candidate pairs are computed, and the pair with the highest predicted score is selected per protein.

\subsubsection{Feature Representation}

Each candidate pair is represented by 25 features across five categories (full listing in Table~S3):

\textbf{A-stage features (8).} Binary presence indicators for each of the four search tools at the locus, plus aggregate metrics: number of tools hitting the locus (\texttt{a\_n\_tools}), best bitscore among all tools (\texttt{a\_best\_score}), best query coverage (\texttt{a\_best\_qcov}), and total raw hit count (\texttt{a\_n\_raw\_hits}). These capture the strength and diversity of evidence from Phase~1.

\textbf{Geometric features (6).} Genomic region size (linear and log-transformed), whether the locus lies on the same chromosome as the primary ortholog, strand agreement with the primary ortholog, bp overlap with the co-ortholog recovery region, and whether the locus shares the same chromosome as the co-ortholog region. These capture spatial priors---for instance, paralogs are often located on different chromosomes or in tandem arrays.

\textbf{Query features (2).} Protein length in amino acids (linear and log-transformed). Longer proteins generate more candidate hits due to having more matchable sequence, and the router must learn to calibrate for this.

\textbf{C-stage features (8).} Spliced alignment quality metrics from the Phase~3 miniprot output, as described in Section~3.2. These are the most informative features: \texttt{c\_identity} is the single most important predictor (mean gain 0.19 across 4 species, Table~S3), followed by \texttt{c\_coverage} (0.12). C-stage features are available for both human-query and species-protein-query runs, enabling the router to assess alignment quality under both conditions.

\textbf{Species indicator (1).} A binary feature indicating whether the query protein used at this locus is the human query or the species-matched protein. This feature allows the router to calibrate its predictions differently for cross-species vs.\ within-species alignments.

\subsubsection{Training Procedure}

\textbf{Training data.} For each protein, we generate candidate pairs from Stage~1 and run miniprot twice per locus---once with the human query protein and once with the species-matched protein. This yields paired (human-query, species-query) feature rows per candidate locus. The training target $y_{L,q}$ is the best-overlap exon F1 of the miniprot spliced alignment at that locus, computed against the ground-truth CDS annotations for the correct paralog gene. The model minimizes mean squared error:
\begin{equation}
\mathcal{L}(\theta) = \frac{1}{|\mathcal{D}|} \sum_{(L, q) \in \mathcal{D}} \bigl( f_\theta(\mathbf{x}_{L, q}) - y_{L, q} \bigr)^2,
\label{eq:mse_loss}
\end{equation}
where $\mathcal{D}$ is the training candidate pair set. Across all four species, the training set contains 25,732 candidate regions from 1,901 proteins (4,511 mouse, 7,910 dog, 5,844 chicken, 7,463 zebrafish).

\textbf{Model.} We use a gradient-boosted tree ensemble (XGBoost) as the regressor. XGBoost is chosen over neural networks for three reasons: (1) it handles tabular features without requiring normalization or extensive hyperparameter tuning; (2) it provides interpretable feature importance scores via total gain (loss reduction from splits on each feature); and (3) it is robust to the small-to-moderate training set sizes in our per-species setting (392--399 training proteins per species, $\sim$3,600--6,300 training rows). The model uses 200 trees with maximum depth 6, learning rate 0.05, subsample 0.8, and column subsample 0.8.

\textbf{Validation.} We use GroupShuffleSplit with an 80/20 split, grouped by protein ID to prevent data leakage: all candidate pairs for a given protein appear exclusively in either the training or test fold. This ensures the router is evaluated on proteins it has never seen, testing its ability to generalize routing decisions across the proteome.

\subsubsection{Inference}

At inference time, for a test protein, we:
\begin{enumerate}
    \item Generate candidate pairs via Stage~1 (merge A-stage tool hits, apply co-ortholog recovery, run dual miniprot queries).
    \item For each candidate pair, extract the 25-dimensional feature vector.
    \item The XGBoost model predicts a scalar exon F1 for each pair.
    \item The router selects the pair with the highest predicted F1, and returns the corresponding miniprot spliced alignment.
\end{enumerate}

The router achieves 98.7\% of oracle F1 on mouse, indicating that the 25 features contain sufficient signal to nearly saturate the routing ceiling. In mouse, 87 of 97 test proteins (89.7\%) match the oracle's exact candidate selection---the remainder differ in their choice but achieve comparable F1. The router substantially outperforms the strongest non-learned heuristic: max-identity locus selection achieves primary F1 of 0.36, while XGBoost routing achieves 0.73 (+0.37 mean across species, Table~S4). This gap demonstrates that the router captures feature interactions beyond what any single heuristic can express---for instance, a locus on the primary chromosome with moderate identity may be preferred over a high-identity locus on a different chromosome, a conditional pattern that only a learned model can represent. This directly addresses the H3 oracle gap identified in Exp~6 (Table~\ref{tab:s_exp6_hybrid}): the per-protein oracle ceiling of $+0.23$--$0.29$ over the best fixed hybrid is substantially captured by learned routing. The router's feature importance analysis and biological interpretation are presented in Section~4.2.

\subsection{Algorithm Summary}

Algorithm~\ref{alg:tprs} presents the complete TPRS procedure.

\begin{algorithm}[t]
\caption{TPRS: Candidate Generation + Learned Joint Selection}
\label{alg:tprs}
\begin{algorithmic}[1]
\small
\Require Human query protein $Q$, target genome $G$, species $S$
\Require A-stage tools $\mathcal{T} = \{\text{diamond}, \text{last}, \text{miniprot}, \text{mmseqs}\}$ (relaxed params)
\Ensure Best spliced alignment and its predicted exon F1
\State \textbf{// Stage 1: Candidate Generation}
\State $\mathcal{R} \gets \emptyset$ \Comment{raw hits}
\For{each tool $t \in \mathcal{T}$}
    \State $\mathcal{R} \gets \mathcal{R} \cup \text{run}_t(Q, G)$ \Comment{relaxed sensitivity}
\EndFor
\State $\mathcal{L} \gets \text{MergeOverlapping}(\mathcal{R}, \text{slop}=50\text{kb})$ \Comment{by chrom+strand}
\If{$\nexists L \in \mathcal{L}$ overlapping primary ortholog locus}
    \State $P_{\text{co}} \gets \text{ComparaCoOrthologProtein}(Q, S)$ \Comment{via proteome}
    \State $\mathcal{L} \gets \mathcal{L} \cup \text{MMseqs2Search}(P_{\text{co}}, G)$
\EndIf
\State $\mathcal{C} \gets \emptyset$ \Comment{candidate pool}
\For{each locus $L \in \mathcal{L}$}
    \State $g \gets \text{BestMatchingGene}(L, \text{CDS\_annotations}_S)$
    \State $P_{\text{sp}} \gets \text{SpeciesProteome}_S(g)$
    \For{each query $q \in \{Q, P_{\text{sp}}\}$}
        \State $\text{align} \gets \text{miniprot}(q, G[L.\text{start}:L.\text{end}], \texttt{-N~3})$
        \State $\mathbf{x}_{(L,q)} \gets \text{ExtractFeatures}(L, \text{align})$ \Comment{25-dim}
        \State $\mathcal{C} \gets \mathcal{C} \cup \{(L, q, \mathbf{x}_{(L,q)}, \text{align})\}$
    \EndFor
\EndFor
\State \textbf{// Stage 2: Learned Joint Selection}
\For{each candidate $(L,q,\mathbf{x},\text{align}) \in \mathcal{C}$}
    \State $\hat{y}_{(L,q)} \gets \text{XGBoost}(\mathbf{x})$ \Comment{predicted exon F1}
\EndFor
\State $(L^*, q^*) \gets \arg\max_{(L,q)} \hat{y}_{(L,q)}$
\State \Return $\text{align}_{(L^*,q^*)}$ and $\hat{y}_{(L^*,q^*)}$
\end{algorithmic}
\end{algorithm}
